%% AI工具使用详情.tex (v1.5.7.3 通用模板)
%% 2026 国赛 AI 工具使用规定第 4 条
%% 配套 skill: li-mtrie-2026 v1.5.7.3
%% 编译: xelatex -interaction=nonstopmode AI工具使用详情.tex (×2)
%% 输出: 支撑材料/AI工具使用详情.pdf (≈ 4-6 页)
%% 模板: 所有具体内容用【】占位符 + \Placeholder{} 宏, 跑题用户复制后替换
%% 评阅要求: 必须真实可查, 不能虚构 AI 用途或夸大 AI 占比, 否则按违规处理
%% 2026 规范第十条: 本文件单独成 PDF 放支撑材料/, 跟 论文/电子版 平行提交
%% preamble-first.tex (v1.5.7.3 — 必须在 \documentclass 之前)
%% \PassOptionsToPackage{quiet}{xeCJK} 是给 ctex 包读的, 必须在 \documentclass 之前
\PassOptionsToPackage{quiet}{xeCJK}


\documentclass[electronic]{format}

%% preamble.tex (v1.5.7.3 — 抽自 论文.tex / 电子版.tex / AI工具使用详情.tex, 在 \documentclass 之后)
%% 三个入口文件共享的 preamble 部分:
%%   - 教学开关 (\Teach / \TeachExample / \Placeholder) + 默认隐藏
%%   - 兜底宏 (*Slot 4 个空定义, 跑题用户忘删不编译失败)
%%   - 公共包 (etoolbox / ctex / mdframed / url / array / tabularx / longtable / tikz)
%%   - 公共列定义 (C / R / L for tabularx) + bold 字体重定义
%% 3 个文件各自保留:
%%   - 自己的 \documentclass[选项]{format} (withpreface / electronic / electronic)
%%   - 自己的注释 (2026 规范说明)
%%   - 自己的 \title{}
%% 编译入口仍是 论文.tex / 电子版.tex / AI工具使用详情.tex, preamble-first.tex + preamble.tex 由 \input{} 拉入

%% === v1.4.4 借鉴自 数模陪跑独家 模板：教学开关 + 占位符宏 ===
%% 在文档任意位置可用 \Teach{...} 写红色填写提示，\TeachExample{...} 写红色示例。
%% \Placeholder{...} 用黑色加粗显示【】占位符，便于交稿前搜索替换。
%% 默认 \showteachingfalse：教学提示不显示（不影响现有 PDF 输出）。
%% 想看教学提示：把下面 \showteachingfalse 注释掉，把上面 \showteachingtrue 取消注释即可。
\newif\ifshowteaching
\showteachingfalse
% \showteachingtrue

\newcommand{\Placeholder}[1]{\textbf{#1}}

\newcommand{\Teach}[1]{%
  \ifshowteaching
    \par\begingroup
      \noindent\color{red!78!black}\bfseries
      【填写提示（提交前自动隐藏）：#1】\par
    \endgroup
  \fi
}

%% === v1.5.7 兜底宏: 跑题用户忘删 *Slot 行不会编译失败 (空定义, 不影响输出) ===
\newcommand{\AISlot}[1]{}
\newcommand{\AppendixSlot}[1]{}
\newcommand{\CodeSlot}[1]{}
\newcommand{\RefSlot}[1]{}

\newcommand{\TeachExample}[1]{%
  \ifshowteaching
    \par\begingroup
      \noindent\color{red!78!black}\bfseries
      【示例（仅供理解，提交前自动隐藏）：#1】\par
    \endgroup
  \fi
}
%% === v1.4.4 借鉴结束 ===

%% === v1.5.7.3 公共包 + 列定义 + tikz (从原 3 文件 preamble 抽出) ===
\usepackage{etoolbox}
\usepackage{ctex}
\BeforeBeginEnvironment{tabular}{\zihao{-5}}
\usepackage[framemethod=TikZ]{mdframed}
\usepackage{url}
\usepackage{array}
\usepackage{tabularx}
\usepackage{longtable}
\newcolumntype{C}{>{\centering\arraybackslash}X}
\newcolumntype{R}{>{\raggedleft\arraybackslash}X}
\newcolumntype{L}{>{\raggedright\arraybackslash}X}

\renewcommand{\textbf}[1]{{\song\bfseries #1}}

\usepackage{tikz}
\usetikzlibrary{arrows.meta,positioning}
\tikzset{
  box/.style={rectangle, draw, minimum width=3.2cm, minimum height=0.9cm, align=center},
  arrow/.style={thick, -{Stealth}}
}
%% === v1.5.7.3 抽出结束 ===


\title{AI 工具使用详情}

\begin{document}

\maketitle
\thispagestyle{plain}
\vspace{-2em}

%% ====== §1 工具信息 (5-8 行表格, 每行可增减) ======
\section{工具信息}

\begin{longtable}{>{\centering\arraybackslash}p{0.25\textwidth}>{\raggedright\arraybackslash}p{0.70\textwidth}}
  \toprule
  项目 & 内容 \\
  \midrule
  \Placeholder{【工具名称, 如 ChatGPT / Claude / DeepSeek / Mavis 等】} & \Placeholder{【填工具官方全名】} \\
  \Placeholder{提供方} & \Placeholder{【填公司/组织全名】} \\
  \Placeholder{版本/型号} & \Placeholder{【填模型版本号】} \\
  \Placeholder{接入方式} & \Placeholder{【填 web/桌面/API/插件 等】} \\
  \Placeholder{使用时间} & \Placeholder{【填 YYYY-MM-DD ~ YYYY-MM-DD (赛题时间窗)】} \\
  \Placeholder{【其他, 如账号/付费情况/上下文长度】} & \Placeholder{【填】} \\
  \bottomrule
\end{longtable}

%% ====== §2 5 项自检 (2026 规范 5 项硬要求) ======
\section{AI 参与边界自检}

%% 借鉴 BZD 数模社 论文正文-AI工具使用声明 2026 规范: 5 项自检, 详细使用情况见支撑材料

\begin{longtable}{>{\raggedright\arraybackslash}p{0.55\textwidth}>{\centering\arraybackslash}p{0.10\textwidth}>{\raggedright\arraybackslash}p{0.28\textwidth}}
  \toprule
  \textbf{自检项} & \textbf{是/否} & \textbf{说明} \\
  \midrule
  AI 是否参与核心模型设计 & \Placeholder{【是/否】} & \Placeholder{【填: 模型思路/公式推导/参数选择是否由本参赛队独立完成; AI 仅辅助代码实现/查资料 = 否; AI 直接给模型 = 是】} \\
  AI 是否生成或修改程序代码 & \Placeholder{【是/否】} & \Placeholder{【填: 详见 §3 使用目的, 列出哪几个 Python 文件 AI 参与】} \\
  AI 输出是否经人工审查、修改与核验 & \Placeholder{【是/否】} & \Placeholder{【填: 详见 §4 核验情况表, 至少 1 轮人工 review + 单元测试/解析解对照】} \\
  AI 是否独立完成核心建模 / 公式推导 / 程序计算 / 最终结论 & \Placeholder{【是/否】} & \Placeholder{【填: 公式推导/参数选择/最终结论均由本参赛队完成; AI 仅辅助 = 否; 关键结论靠 AI = 是】} \\
  AI 使用占比 (估算) & -- & \Placeholder{【填: 资料检索 X\% + 代码编写 Y\% + 文档润色 Z\% + 其他 W\% = 100\%; 最终决策与判定 100\% 由参赛队完成】} \\
  \bottomrule
\end{longtable}

%% ====== §3 三阶段使用目的 (BZD 数模社规范: 资料与思路 → 数据与程序 → 检验与写作) ======
\section{使用目的与环节}

本参赛队在以下环节使用了 AI 工具（按 BZD 规范分为"资料与思路 → 数据与程序 → 检验与写作"三阶段）。\textbf{每一项条目都必须对应 §4 中一条具体 prompt}，不能空泛列项。

\textbf{第一阶段：资料与思路}
\begin{itemize}[itemindent=2em]
  \item \textbf{\Placeholder{【环节名, 如: 题目理解】}}: \Placeholder{【填: AI 做了什么, 如辅助读题 (PyMuPDF 提取 PDF 文本)/ 识别问题数 / 判定题组与题号】}
  \item \textbf{\Placeholder{【环节名, 如: 模型选择】}}: \Placeholder{【填: AI 做了什么, 如候选模型对比表初稿 / 文献综述 / 思路启发】}
  \item \Placeholder{【可增删更多环节, 如: 文献检索 / 题意翻译 / 创新点挖掘】}
\end{itemize}

\textbf{第二阶段：数据与程序}
\begin{itemize}[itemindent=2em]
  \item \textbf{\Placeholder{【环节名, 如: 代码编写】}}: \Placeholder{【填: AI 写了哪几个 .py / 哪个函数 / 哪段算法, 如数据预处理 / 优化求解 / 灵敏度分析】}
  \item \textbf{\Placeholder{【环节名, 如: 数据预处理】}}: \Placeholder{【填: AI 做了什么, 如 profile\_data.py 列类型识别 / 缺失值填补 / 异常值检测】}
  \item \textbf{\Placeholder{【环节名, 如: 结果可视化】}}: \Placeholder{【填: AI 画了哪些图, 调整了什么参数 (DPI/字号/线宽)】}
  \item \Placeholder{【可增删更多环节, 如: 公式排版 / 表格生成 / 参考文献格式化】}
\end{itemize}

\textbf{第三阶段：检验与写作}
\begin{itemize}[itemindent=2em]
  \item \textbf{\Placeholder{【环节名, 如: 文档润色】}}: \Placeholder{【填: AI 做了什么, 如 LaTeX 公式排版 / longtable 格式 / GB/T 7714 参考文献】}
  \item \textbf{\Placeholder{【环节名, 如: 查重改写】}}: \Placeholder{【填: AI 做了什么, 如问题重述改写 / 避免与赛题原文重复 / AIGC 痕迹降低】}
  \item \textbf{\Placeholder{【环节名, 如: 合规检查】}}: \Placeholder{【填: AI 做了什么, 如按 references/合规检查清单.md 4 类硬规则逐项核对】}
  \item \Placeholder{【可增删更多环节, 如: 摘要改写 / 答辩准备 / 5 步状态机 smoke test】}
\end{itemize}

%% ====== §4 主要 prompt 与过程 (节选) ======
\section{主要 prompt 与过程 (节选)}

%% 评阅只看 §4 验证 §3 真实性. §3 列的每条环节, §4 必须有对应 prompt 痕迹.
%% 一条 prompt = 一个 \subsection. 删掉不用的, 复制现有格式加新的.

\subsection{3.1 \Placeholder{【环节类型】} prompt: \Placeholder{【任务简述】}}

\textbf{Prompt}: \Placeholder{【把发给 AI 的原始 prompt 完整贴这里, 1-3 行为佳; 不要改写, 评阅会查真实性】}

\textbf{AI 响应}:
\begin{itemize}
  \item \Placeholder{【AI 给了什么, 第 1 点】}
  \item \Placeholder{【AI 给了什么, 第 2 点】}
  \item \Placeholder{【AI 给了什么, 第 3 点】}
\end{itemize}

\textbf{采纳 / 修改}:
\begin{itemize}
  \item \textbf{采纳}: \Placeholder{【哪些点直接用】}
  \item \textbf{关键修改}: \Placeholder{【哪些点改了, 改了之后怎样; 改 AI 代码 bug / 调整算法参数 / 修正公式 等】}
  \item \textbf{效果}: \Placeholder{【改前 X / 改后 Y / 提升 Z\%; 用数据说话】}
\end{itemize}

\subsection{3.2 \Placeholder{【环节类型】} prompt: \Placeholder{【任务简述】}}

\textbf{Prompt}: \Placeholder{【贴原始 prompt】}

\textbf{AI 响应}: \Placeholder{【贴 AI 主要回答, 1-3 行为佳; 长回答可摘关键句】}

\textbf{采纳}: \Placeholder{$\checkmark$ / $\times$ / 部分} \Placeholder{【填: 全部采纳 / 部分采纳 / 不用; 给出 1-2 句理由】}

\Placeholder{【如果此 prompt 直接产生了可量化的结果 (如优化提升/精度提升), 在此加一段:】}
\textbf{\Placeholder{【结果】}}:
\begin{itemize}
  \item \Placeholder{【Baseline: 数值】}
  \item \Placeholder{【优化后: 数值】} (\Placeholder{【+/-X\%】})
\end{itemize}

%% 如还有更多 prompt, 复制 3.1 / 3.2 格式继续加 3.3 / 3.4 / ...

%% ====== §5 核验情况表 (与 §3 环节对应, 评阅重点查真实性) ======
\section{对 AI 输出的核验情况}

所有 AI 生成的内容均经过人工核验和修改, 具体如下。\textbf{每一条核验条目必须能在 §3 / §4 找到对应来源}。

\begin{longtable}{>{\raggedright\arraybackslash}p{0.32\textwidth}>{\raggedright\arraybackslash}p{0.63\textwidth}}
  \toprule
  \textbf{AI 输出} & \textbf{核验方式与修改内容} \\
  \midrule
  \Placeholder{【AI 输出 1, 对应 §3 第 X 阶段第 Y 条】} & \Placeholder{【填: 单元测试 + 解析解对照 / 多次运行检查 / 人工审阅 + 规范核对; 改了哪里, 为什么改】} \\
  \Placeholder{【AI 输出 2】} & \Placeholder{【填核验方式】} \\
  \Placeholder{【AI 输出 3】} & \Placeholder{【填核验方式】} \\
  \Placeholder{【AI 输出 4】} & \Placeholder{【填核验方式】} \\
  \Placeholder{【AI 输出 5】} & \Placeholder{【填核验方式】} \\
  \Placeholder{【AI 输出 6】} & \Placeholder{【填核验方式】} \\
  \Placeholder{【可增删更多核验条目, 每条对应 §3 一个环节】} & \Placeholder{【填】} \\
  \bottomrule
\end{longtable}

\textbf{核验流程} (本参赛队通用流程, 所有 AI 输出必经):
\begin{enumerate}
  \item AI 给出代码 / 公式后, \textbf{手算} 1-2 个简单 case 验证 (边界条件 + 典型场景)
  \item 跑 \texttt{py} smoke test, 跑 \texttt{py tools/pack.py} 验证 zip
  \item 跑 \texttt{grep '! Error|Overfull'} 查 LaTeX 编译日志
  \item 与队友 / 教练讨论, 不通过就改
\end{enumerate}

\vspace{2em}

\noindent\rule{\textwidth}{0.5pt}

\noindent\textit{参赛队号: \rule{3cm}{0.4pt}\hfill 日期: \rule{2cm}{0.4pt}}

\noindent\textit{参赛队员: \rule{3cm}{0.4pt} \rule{3cm}{0.4pt} \rule{3cm}{0.4pt} \hfill 指导教师: \rule{2cm}{0.4pt}}

\end{document}
